
\documentclass{article}
\usepackage{colortbl}
\usepackage{makecell}
\usepackage{multirow}
\usepackage{supertabular}

\begin{document}

\newcounter{utterance}

\centering \large Interaction Transcript for game `cladder', experiment `full\_v1.5\_default', episode 4371 with refineonly{-}Qwen3.5{-}27B.
\vspace{24pt}

{ \footnotesize  \setcounter{utterance}{1}
\setlength{\tabcolsep}{0pt}
\begin{supertabular}{c@{$\;$}|p{.15\linewidth}@{}p{.15\linewidth}p{.15\linewidth}p{.15\linewidth}p{.15\linewidth}p{.15\linewidth}}
    \# & \multicolumn{2}{c}{Player} && \multicolumn{2}{c}{Game Master} \\
    \hline

    \theutterance \stepcounter{utterance}  
    & & & \multicolumn{4}{p{0.6\linewidth}}{
        \cellcolor[rgb]{0.9,0.9,0.9}{
            \makecell[{{p{\linewidth}}}]{
                \texttt{\tiny{[P1$\langle$GM]}}
                \texttt{You are an expert in causal inference. The following question is not a typical commonsense query, but rather a meticulously designed question created by a professor specializing in causal inference, intended to assess the students' mastery of the course content.} \\
\texttt{Imagine a self{-}contained, hypothetical world with only the following conditions, and without any unmentioned factors or causal relationships: Gender has a direct effect on smoking and peanut allergy. Smoking has a direct effect on tar deposit. Tar deposit has a direct effect on peanut allergy. Gender is unobserved. For nonsmokers, the probability of high tar deposit is 28\%. For smokers, the probability of high tar deposit is 58\%. For nonsmokers and with no tar deposit, the probability of being allergic to peanuts is 49\%. For nonsmokers and with high tar deposit, the probability of being allergic to peanuts is 77\%. For smokers and with no tar deposit, the probability of being allergic to peanuts is 27\%. For smokers and with high tar deposit, the probability of being allergic to peanuts is 62\%. The overall probability of smoking is 58\%. Will smoking increase the chance of being allergic to peanuts?} \\
\texttt{Start your answer with 'Yes' or 'No', followed by additional reasoning or evidence to support your explanation.} \\
            }
        }
    }
     \\ \\

    \theutterance \stepcounter{utterance}  
    & \multicolumn{4}{p{0.6\linewidth}}{
        \cellcolor[rgb]{0.561,0.737,0.561}{
            \makecell[{{p{\linewidth}}}]{
                \texttt{\tiny{[P1$\rangle$GM]}}
                \texttt{No} \\
\\ 
\texttt{To determine whether smoking increases the chance of being allergic to peanuts, we must calculate the **Average Treatment Effect (ATE)** of smoking on peanut allergy. This requires comparing the expected probability of allergy under the intervention of smoking (\$P(Y=1|do(S=1))\$) versus the intervention of not smoking (\$P(Y=1|do(S=0))\$).} \\
\\ 
\texttt{**1. Analyze the Causal Structure**} \\
\texttt{The problem describes the following causal graph:} \\
\texttt{*   \$Gender \textbackslash{}to Smoking\$} \\
\texttt{*   \$Gender \textbackslash{}to PeanutAllergy\$} \\
\texttt{*   \$Smoking \textbackslash{}to TarDeposit\$} \\
\texttt{*   \$TarDeposit \textbackslash{}to PeanutAllergy\$} \\
\\ 
\texttt{Here, \$Gender\$ is an unobserved confounder between \$Smoking\$ and \$PeanutAllergy\$. The back{-}door path is \$S \textbackslash{}leftarrow Gender \textbackslash{}to Y\$. Since \$Gender\$ is unobserved, we cannot condition on it directly to block this path. However, we can use the **Front{-}Door Criterion** because \$TarDeposit\$ (\$T\$) acts as a mediator that intercepts the causal path from \$S\$ to \$Y\$.} \\
\\ 
\texttt{The conditions for the Front{-}Door Criterion are met:} \\
\texttt{1.  \$T\$ blocks all directed paths from \$S\$ to \$Y\$ (since the only path is \$S \textbackslash{}to T \textbackslash{}to Y\$).} \\
\texttt{2.  There are no unblocked back{-}door paths from \$S\$ to \$T\$ (the only path is \$S \textbackslash{}leftarrow Gender \textbackslash{}to Y \textbackslash{}leftarrow T\$, which is blocked by the collider \$Y\$? No, wait. The path is \$S \textbackslash{}leftarrow Gender \textbackslash{}to Y \textbackslash{}leftarrow T\$. This is not a backdoor path from S to T. The backdoor path from S to T would be \$S \textbackslash{}leftarrow Gender \textbackslash{}dots \textbackslash{}to T\$. There is no arrow from Gender to T mentioned. The text says "Gender has a direct effect on smoking and peanut allergy." It does not mention an effect on Tar Deposit. Thus, \$S\$ and \$T\$ are not confounded by \$Gender\$).} \\
\texttt{3.  All back{-}door paths from \$T\$ to \$Y\$ are blocked by \$S\$ (the path \$T \textbackslash{}leftarrow S \textbackslash{}leftarrow Gender \textbackslash{}to Y\$ is blocked by conditioning on \$S\$).} \\
\\ 
\texttt{Therefore, the causal effect is identified using the Front{-}Door formula:} \\
\texttt{\$\$P(Y|do(S)) = \textbackslash{}sum\_\{t\} P(Y|S, t) P(t|S)\$\$} \\
\texttt{*Note: Since \$T\$ is the only mediator and there is no direct edge \$S \textbackslash{}to Y\$ mentioned in the "direct effect" list (only \$S \textbackslash{}to T\$ and \$T \textbackslash{}to Y\$), the effect of \$S\$ on \$Y\$ is entirely mediated by \$T\$. Even} \\
            }
        }
    }
    & & \\ \\

    \theutterance \stepcounter{utterance}  
    & & & \multicolumn{2}{p{0.3\linewidth}}{
        \cellcolor[rgb]{0.9,0.9,0.9}{
            \makecell[{{p{\linewidth}}}]{
                \texttt{\tiny{[GM$|$GM]}}
                \texttt{No} \\
            }
        }
    }
    & & \\ \\

    \theutterance \stepcounter{utterance}  
    & & & \multicolumn{2}{p{0.3\linewidth}}{
        \cellcolor[rgb]{0.9,0.9,0.9}{
            \makecell[{{p{\linewidth}}}]{
                \texttt{\tiny{[GM$|$GM]}}
                \texttt{yes} \\
            }
        }
    }
    & & \\ \\

    \theutterance \stepcounter{utterance}  
    & & & \multicolumn{2}{p{0.3\linewidth}}{
        \cellcolor[rgb]{0.9,0.9,0.9}{
            \makecell[{{p{\linewidth}}}]{
                \texttt{\tiny{[GM$|$GM]}}
                \texttt{game\_result = LOSE} \\
            }
        }
    }
    & & \\ \\

\end{supertabular}
}

\end{document}
